\section{Attribution thành phần}
\label{sec:ch07-attribution-thanh-phan}

\index{attribution thành phần}%
\index{vá kích hoạt}%
\index{diễn giải cơ chế}%
\tn{Attribution dữ liệu}{data attribution} truy về giai đoạn huấn luyện, còn
\tn{attribution thành phần}{component attribution} ở lại trong một lượt suy diễn
và nhìn vào bên trong mô hình. Đối tượng ở đây là các bộ phận của chính mạng: một nơ-ron, một
tầng, một đầu chú ý, hoặc một mạch, tức một mạng con nối các bộ phận ấy lại thành
một đường xử lý~\autocite{p18unified}. Đây là nhánh mà
\tn{diễn giải cơ chế}{mechanistic interpretability} coi là trung tâm, vì nó hứa
một câu trả lời về cơ chế chứ không chỉ về mối liên hệ giữa đầu vào và đầu ra.

Phép nhiễu ở nhánh này mang tên
\tn{phân tích trung gian nhân quả}{causal mediation analysis}, và nó chạy trên
ba lượt. Lượt thứ nhất là lượt sạch, không can thiệp gì. Lượt thứ hai làm hỏng
các thành phần, có thể bằng cách làm hỏng đầu vào để sinh ra kích hoạt sai. Lượt
thứ ba giữ nguyên tình huống đã hỏng nhưng chép giá trị kích hoạt của một thành
phần từ lượt sạch trở lại, rồi đo xem đầu ra hồi phục được bao
nhiêu~\autocite{p18unified}. Thao tác ở lượt thứ ba được gọi là
\tn{vá kích hoạt}{activation patching}; khi nó áp lên các đường nối giữa nhiều
thành phần thay vì lên một thành phần đơn lẻ thì tên gọi là
\tn{vá theo đường}{path patching}.

Điểm yếu của thủ tục nằm ở chữ \enquote{làm hỏng}. Bài báo liệt kê ít nhất năm cách: đặt
kích hoạt về không, thay bằng trung bình, thêm nhiễu Gauss đã làm trơn, tráo
kích hoạt từ một đầu vào khác, hoặc học một phép cắt bỏ. Bộ dữ liệu dùng để sinh
kích hoạt cũng phải được chọn sao cho nó thật sự gợi ra hành vi cần xét, kèm một
chỉ số đo hành vi ấy~\autocite{p18unified}. Không có cách nào trong số đó tự
chứng minh mình là cách đúng, nên lựa chọn lại rơi vào tay người làm thí nghiệm,
đúng như baseline trong chương~\ref{ch:attribution-theo-gradient}
và độ đo trong chương~\ref{ch:attribution-tu-nguyen-ly-dau}.

Hai kỹ thuật còn lại cũng có mặt. Neuron Shapley mang giá trị Shapley vào nhánh
này, với phép lấy mẫu và một thuật toán multi-armed bandit, tức thủ tục dồn
dần phép lấy mẫu vào những nơ-ron đang có điểm cao thay vì lấy mẫu đều, để
khoanh nhanh những nơ-ron điểm cao. Về phía gradient, attribution patching xấp
xỉ thay đổi cục bộ của một lần vá thay vì chạy lại mô hình cho từng thành
phần: nó lấy gradient của chỉ số vá theo các kích hoạt ở lượt bị hỏng, lưu
lại, rồi nhân với hiệu giữa kích hoạt sạch và kích hoạt hỏng của từng thành
phần, nên chỉ cần hai lượt truyền xuôi cho hai bộ kích hoạt và một lượt truyền
ngược cho gradient; bản mở rộng cho mọi cạnh của đồ thị tính toán
chấm điểm tất cả
rồi giữ lại $k$ cạnh cao nhất làm mạch~\autocite{p18unified}. Về phía
\tn{xấp xỉ tuyến tính}{linear approximation}, COAR khớp một mô hình tuyến tính
dự đoán tác động phản thực của việc cắt bỏ từng thành phần, rồi cộng điểm số để
ước lượng tác động của cả một nhóm.

Một ranh giới trong bài báo đáng ghi lại, vì Phần~V sẽ quay lại nó. Sparse
autoencoder được huấn luyện để dựng lại kích hoạt của mô hình dưới ràng buộc
thưa, và các tác giả không xếp nó vào attribution thành phần: nó học ra thành
phần mới thay vì gán điểm cho thành phần đã có. Vai trò của nó là tìm ra những
thành phần dễ diễn giải, sau đó mới đưa vào một phương pháp
attribution~\autocite{p18unified}. Chương~\ref{ch:nut-that-khai-niem} đọc kỹ họ
mô hình này khi bàn về nút thắt khái niệm.

Đến đây ba nhánh đã có đủ ví dụ cho cả ba kỹ thuật. Nhưng lặp lại kỹ thuật mới
là quan sát, chưa phải một khung, nên bước tiếp theo là khung hình thức duy nhất
mà bài báo viện đến.
